\begin{figure}[H]
    \centering
    \begin{tikzpicture}[x=1.0cm, y=1.0cm]

        % input
        \foreach \v [count=\i from 0] in {5,2,9,1,7,3,8,4} {
            \node[lane] (in\i) at (0.76*\i,0) {$\v$};
        }
        \node[rowlab] at (in0.west) {είσοδος};

        % (a) full sort
        \node[panel, font=\scriptsize\bfseries, anchor=west] at (-1.05,-1.42)
            {(α) ταξινόμηση};
        \node[notemute, font=\tiny, anchor=east] at (5.68,-1.42)
            {$\Omega(n \log n)$ συγκρίσεις};
        \foreach \v [count=\i from 0] in {9,8,7,5,4,3,2,1} {
            \node[lanereg] (so\i) at (0.76*\i,-1.90) {$\v$};
        }
        \draw[bracemirror] (-0.40,-2.38) -- (5.72,-2.38);
        \node[notemute, anchor=north, font=\tiny] at (2.66,-2.62)
            {και οι $n$ θέσεις προσδιορισμένες};

        % (b) selection
        \node[panel, font=\scriptsize\bfseries, anchor=west] at (-1.05,-3.26)
            {(β) επιλογή \tl{top-k}, $k=3$};
        \node[notemute, font=\tiny, anchor=east] at (5.68,-3.26)
            {$\Theta(n)$ συγκρίσεις για σταθερό \tl{k}};
        \foreach \v [count=\i from 0] in {9,7,8} {
            \node[lanereg] (se\i) at (0.76*\i,-3.74) {$\v$};
        }
        \foreach \i in {3,4,5,6,7} {
            \node[lanedrop] at (0.76*\i,-3.74) {$\cdot$};
        }
        \draw[bracemirror] (-0.40,-4.22) -- (1.90,-4.22);
        \node[notemute, anchor=north, font=\tiny, align=center] at (0.76,-4.46)
            {τα \tl{k} μεγαλύτερα,\\ σε οποιαδήποτε σειρά};
        \draw[bracemirror] (1.90,-4.22) -- (5.72,-4.22);
        \node[notemute, anchor=north, font=\tiny, align=center] at (3.81,-4.46)
            {$n-k$ στοιχεία, καμία\\ πληροφορία σειράς};

    \end{tikzpicture}
    \caption{Η επιλογή ζητά αυστηρά λιγότερη πληροφορία από την ταξινόμηση. Από
    την ίδια είσοδο, η ταξινόμηση (α) προσδιορίζει τη σχετική σειρά και των $n$
    στοιχείων, ενώ η επιλογή (β) αρκείται στον διαχωρισμό τους σε δύο ομάδες: τα
    \tl{k} μεγαλύτερα, χωρίς καμία απαίτηση για τη σειρά τους μεταξύ τους, και τα
    υπόλοιπα, για τα οποία δεν χρειάζεται να γνωρίζουμε τίποτε. Η διαφορά αυτή
    αποτυπώνεται στα κάτω φράγματα και εξηγεί γιατί κάθε υλοποίηση που ταξινομεί
    ολόκληρη την είσοδο επιτελεί περισσότερη εργασία από όση απαιτεί το
    πρόβλημα.}
    \label{fig:bg-selection}
\end{figure}
